% Source: https://artofproblemsolving.com/community/c3h1761117_more_number_theory_help
% Stable identity: aops-topic-1761117-q01
Each multiple of $5$ up to $201$ contributes at least one factor $5$; there are $\left\lfloor\dfrac{201}{5}\right\rfloor=40$ of them. But multiples of $25$ contain $5^2$ and contribute a second factor each: there are $\left\lfloor\dfrac{201}{25}\right\rfloor=8$ of them. Multiples of $125$ contribute a third factor: $\left\lfloor\dfrac{201}{125}\right\rfloor=1$ (namely $125$). There are no multiples of $625$. So the exponent is
\[
40+8+1=\boxed{49}.
\]
For $10!$ the only multiples of $5$ are $5$ and $10$, each containing a single $5$; no number up to $10$ is a multiple of $25$, so the extra terms $\left\lfloor\dfrac{10}{25}\right\rfloor,\dots$ are $0$ and $\left\lfloor\dfrac{10}{5}\right\rfloor=2$ is already the full answer. In general the exponent of a prime $p$ in $N!$ is $\left\lfloor\dfrac {N}{p}\right\rfloor+\left\lfloor\dfrac {N}{p^2}\right\rfloor+\left\lfloor\dfrac {N}{p^3}\right\rfloor+\cdots$ (Legendre's formula).
